\documentclass{amsart}
% For dealing with references we use the comment environment
\usepackage{verbatim}
\newenvironment{reference}{\comment}{\endcomment}
%\newenvironment{reference}{}{}
\newenvironment{slogan}{\comment}{\endcomment}
\newenvironment{history}{\comment}{\endcomment}

% For commutative diagrams we use Xy-pic
\usepackage[all]{xy}

% We use 2cell for 2-commutative diagrams.
\xyoption{2cell}
\UseAllTwocells

% We use multicol for the list of chapters between chapters
\usepackage{multicol}

% This is generally recommended for better output
\usepackage{lmodern}
\usepackage[T1]{fontenc}

% For cross-file-references

% Package for hypertext links:
\usepackage{hyperref}

% For any local file, say "hello.tex" you want to link to please

% Theorem environments.
%
\theoremstyle{plain}
\newtheorem{theorem}[subsection]{Theorem}
\newtheorem{proposition}[subsection]{Proposition}
\newtheorem{lemma}[subsection]{Lemma}

\theoremstyle{definition}
\newtheorem{definition}[subsection]{Definition}
\newtheorem{example}[subsection]{Example}
\newtheorem{exercise}[subsection]{Exercise}
\newtheorem{situation}[subsection]{Situation}

\theoremstyle{remark}
\newtheorem{remark}[subsection]{Remark}
\newtheorem{remarks}[subsection]{Remarks}

\numberwithin{equation}{subsection}

% Macros
%
\def\lim{\mathop{\mathrm{lim}}\nolimits}
\def\colim{\mathop{\mathrm{colim}}\nolimits}
\def\Spec{\mathop{\mathrm{Spec}}}
\def\Hom{\mathop{\mathrm{Hom}}\nolimits}
\def\Ext{\mathop{\mathrm{Ext}}\nolimits}
\def\SheafHom{\mathop{\mathcal{H}\!\mathit{om}}\nolimits}
\def\SheafExt{\mathop{\mathcal{E}\!\mathit{xt}}\nolimits}
\def\Sch{\mathit{Sch}}
\def\Mor{\mathop{\mathrm{Mor}}\nolimits}
\def\Ob{\mathop{\mathrm{Ob}}\nolimits}
\def\Sh{\mathop{\mathit{Sh}}\nolimits}
\def\NL{\mathop{N\!L}\nolimits}
\def\CH{\mathop{\mathrm{CH}}\nolimits}
\def\proetale{{pro\text{-}\acute{e}tale}}
\def\etale{{\acute{e}tale}}
\def\QCoh{\mathit{QCoh}}
\def\Ker{\mathop{\mathrm{Ker}}}
\def\Im{\mathop{\mathrm{Im}}}
\def\Coker{\mathop{\mathrm{Coker}}}
\def\Coim{\mathop{\mathrm{Coim}}}

% Boxtimes
%
\DeclareMathSymbol{\boxtimes}{\mathbin}{AMSa}{"02}

%
% Macros for moduli stacks/spaces
%
\def\QCohstack{\mathcal{QC}\!\mathit{oh}}
\def\Cohstack{\mathcal{C}\!\mathit{oh}}
\def\Spacesstack{\mathcal{S}\!\mathit{paces}}
\def\Quotfunctor{\mathrm{Quot}}
\def\Hilbfunctor{\mathrm{Hilb}}
\def\Curvesstack{\mathcal{C}\!\mathit{urves}}
\def\Polarizedstack{\mathcal{P}\!\mathit{olarized}}
\def\Complexesstack{\mathcal{C}\!\mathit{omplexes}}
% \Pic is the operator that assigns to X its picard group, usage \Pic(X)
% \Picardstack_{X/B} denotes the Picard stack of X over B
% \Picardfunctor_{X/B} denotes the Picard functor of X over B
\def\Pic{\mathop{\mathrm{Pic}}\nolimits}
\def\Picardstack{\mathcal{P}\!\mathit{ic}}
\def\Picardfunctor{\mathrm{Pic}}
\def\Deformationcategory{\mathcal{D}\!\mathit{ef}}

\setlength{\emergencystretch}{3em}
\begin{document}
\title[Stacks Project, Tag 0DJN]{385 --- Pseudo-coherent complexes are controlled homotopy colimits of perfect approximations}
\maketitle
\noindent Copyright (C) 2005--2025 Johan de Jong. Stacks Project authors. Source: \href{https://stacks.math.columbia.edu/tag/0DJN}{Tag 0DJN}.

\noindent Editorial difficulty note (not author text). Difficulty H2: Independent approximations do not automatically form a compatible directed system. The proof uses degree-specific Ext vanishing to factor each previously chosen perfect approximation through the next, then verifies truncation stabilization and the homotopy-colimit comparison..

\noindent Editorial provenance note (not author text). History: excerpt from revision \texttt{a0444\allowbreak 6e57e\allowbreak c1fbc\allowbreak 252a8\allowbreak 71afc\allowbreak ec775\allowbreak 2fb28\allowbreak 07b14}, perfect.tex, lines 4596--4655. Reference presentation was adapted; original-author context and core support blocks were retained or added. The mathematical wording is unchanged. Licensed under GNU FDL 1.2 or later; no invariant sections or cover texts. See ../licenses/COPYING. Context blocks reproduce author definitions and supporting results; they are not additional counted problems.

\begin{definition}
\label{cohomology-definition-pseudo-coherent}
Let $(X, \mathcal{O}_X)$ be a ringed space. Let $\mathcal{E}^\bullet$
be a complex of $\mathcal{O}_X$-modules. Let $m \in \mathbf{Z}$.
\begin{enumerate}
\item We say $\mathcal{E}^\bullet$ is {\it $m$-pseudo-coherent}
if there exists an open covering $X = \bigcup U_i$ and for each $i$
a morphism of complexes
$\alpha_i : \mathcal{E}_i^\bullet \to \mathcal{E}^\bullet|_{U_i}$
where $\mathcal{E}_i^\bullet$ is strictly perfect on $U_i$ and
$H^j(\alpha_i)$ is an isomorphism for $j > m$ and $H^m(\alpha_i)$
is surjective.
\item We say $\mathcal{E}^\bullet$ is {\it pseudo-coherent}
if it is $m$-pseudo-coherent for all $m$.
\item We say an object $E$ of $D(\mathcal{O}_X)$ is
{\it $m$-pseudo-coherent} (resp.\ {\it pseudo-coherent})
if and only if it can be represented by a $m$-pseudo-coherent
(resp.\ pseudo-coherent) complex of $\mathcal{O}_X$-modules.
\end{enumerate}
\end{definition}

\begin{definition}
\label{definition-approximation-holds}
Let $X$ be a scheme. Consider triples $(T, E, m)$ where
\begin{enumerate}
\item $T \subset X$ is a closed subset,
\item $E$ is an object of $D_\QCoh(\mathcal{O}_X)$, and
\item $m \in \mathbf{Z}$.
\end{enumerate}
We say {\it approximation holds for the triple} $(T, E, m)$ if
there exists a perfect object $P$ of $D(\mathcal{O}_X)$ supported on $T$
and a map $\alpha : P \to E$ which induces isomorphisms $H^i(P) \to H^i(E)$
for $i > m$ and a surjection $H^m(P) \to H^m(E)$.
\end{definition}

\begin{definition}
\label{definition-approximation}
Let $X$ be a scheme. We say {\it approximation by perfect complexes holds}
on $X$ if for any closed subset $T \subset X$ with $X \setminus T$
retro-compact in $X$ there exists an integer $r$ such that
for every triple $(T, E, m)$ as in
Definition \ref{definition-approximation-holds} with
\begin{enumerate}
\item $E$ is $(m - r)$-pseudo-coherent, and
\item $H^i(E)$ is supported on $T$ for $i \geq m - r$
\end{enumerate}
approximation holds.
\end{definition}

\begin{lemma}
\label{lemma-pseudo-coherent-hocolim}
Let $X$ be a quasi-compact and quasi-separated scheme.
Let $K \in D(\mathcal{O}_X)$. The following are equivalent
\begin{enumerate}
\item $K$ is pseudo-coherent, and
\item $K = \text{hocolim} K_n$ where
$K_n$ is perfect and $\tau_{\geq -n}K_n \to \tau_{\geq -n}K$
is an isomorphism for all $n$.
\end{enumerate}
\end{lemma}

\begin{proof}
The implication (2) $\Rightarrow$ (1) is true on any ringed space.
Namely, assume (2) holds. Recall that a perfect object of the derived
category is pseudo-coherent, see
Cohomology, Lemma \href{https://stacks.math.columbia.edu/tag/08CQ}{lemma perfect (Tag 08CQ)}.
Then it follows from the definitions that
$\tau_{\geq -n}K_n$ is $(-n + 1)$-pseudo-coherent
and hence $\tau_{\geq -n}K$ is $(-n + 1)$-pseudo-coherent,
hence $K$ is $(-n + 1)$-pseudo-coherent. This is true for
all $n$, hence $K$ is pseudo-coherent, see
Cohomology, Definition \ref{cohomology-definition-pseudo-coherent}.

\medskip\noindent
Assume (1). We start by choosing an approximation
$K_1 \to K$ of $(X, K, -2)$ by a perfect complex $K_1$, see
Definitions \ref{definition-approximation-holds} and
\ref{definition-approximation} and
Theorem \href{https://stacks.math.columbia.edu/tag/08ES}{theorem approximation (Tag 08ES)}.
Suppose by induction we have
$$
K_1 \to K_2 \to \ldots \to K_n \to K
$$
with $K_i$ perfect such that
such that $\tau_{\geq -i}K_i \to \tau_{\geq -i}K$ is an isomorphism
for all $1 \leq i \leq n$. Then we pick $a \leq b$ as in
Lemma \href{https://stacks.math.columbia.edu/tag/09M4}{lemma ext from perfect into bounded QCoh (Tag 09M4)}
for the perfect object $K_n$. Choose an approximation
$K_{n + 1} \to K$ of $(X, K, \min(a - 1, -n - 1))$.
Choose a distinguished triangle
$$
K_{n + 1} \to K \to C \to K_{n + 1}[1]
$$
Then we see that $C \in D_\QCoh(\mathcal{O}_X)$ has
$H^i(C) = 0$ for $i \geq a$. Thus by our choice of $a, b$
we see that $\Hom_{D(\mathcal{O}_X)}(K_n, C) = 0$.
Hence the composition $K_n \to K \to C$ is zero. Hence by
Derived Categories, Lemma \href{https://stacks.math.columbia.edu/tag/0149}{lemma representable homological (Tag 0149)}
we can factor $K_n \to K$ through $K_{n + 1}$
proving the induction step.

\medskip\noindent
We still have to prove that $K = \text{hocolim} K_n$.
This follows by an application of
Derived Categories, Lemma \href{https://stacks.math.columbia.edu/tag/0CRK}{lemma cohomology of hocolim (Tag 0CRK)}
to the functors
$H^i( - ) : D(\mathcal{O}_X) \to \textit{Mod}(\mathcal{O}_X)$
and our choice of $K_n$.
\end{proof}
\end{document}
